\subsubsection{Single-muon selections}
Muons used in this analysis are required to pass the following selections:

\begin{itemize}
  \item Combined muon
  \item Tight muons:
        For details about tight muon selections and performance,
          see Ref.~\cite{PERF-2015-10}.
        For cross-checks and systematic uncertainties, muons passing the
          ``Medium'' requirement are used.
        The plot captions are always explicitly labelled for what quality muons are used.
  \item $\pT > 4~\GeV$
  \item $ |\eta| < 2.4$
  \item A very loose requirement on the transverse impact parameter of $d_0<5$~mm,  as done in previous analyses~\cite{HION-2015-06}.
\end{itemize}




\subsubsection{Muon-pair selections\label{sec:muon_pair_selections}}
Dimuon pairs are formed with both muons meeting the above criteria.
The muon pair is also required to match the triggers listed in 
  Section~\ref{sec:data} using the 
  \href{https://twiki.cern.ch/twiki/bin/view/Atlas/XAODMatchingTool}{TriggerMatchingTool} 
  package with a $d\mathrm{R}$ requirement of 0.1.
Additionally the average \pt\ of the muons in the pair, 
  $\ptbar=(\pt^{\mu_1}+\pt^{\mu_2})/2$ is required to be larger than 5~\GeV.
This \ptbar\ requirement improves the signal to noise ratio, since 
  the pairs for $4<\ptbar<5$~\GeV\ are dominated by combinatorial
  background pairs which will be shown later.
In events with multiple muon-pairs in the event (that pass all selections),
  all pairs are used in the analysis.
For events with muon-pairs that passes all selections,
  approximately 2\% have more than one muon-pair.


The objective of the analysis is to study the $\Dphi=\phia-\phib$
  distribution between muons where $a$ and $b$ label the two muons 
  in the pair.
\iftoggle{IncludeYields}{
  The goal is to quantify how the width and number of pairs on the 
    away-side ($\Dphi=\pi$), varies with \pt\ and centrality.
}
{
  The goal is to quantify how the width of the pair distribution on the 
    away-side ($\Dphi=\pi$), varies with \pt\ and centrality.
}
Figure~\ref{fig:example_2pc} shows example \Dphi\ distributions,
  for \ptab>4~\GeV.
The pairs are also required to have |\Deta|>0.8, as well as 
  few additional cuts that are discussed later 
  in this section\footnote{
    Typically two-particle correlation analyses use an event-mixing technique
      to account for detector acceptance effects in \Dphi. 
    However the full azimuthal acceptance of ATLAS makes these 
      effects negligible, as has been demonstrated in prior flow 
      analyses in Pb+Pb collisions~\cite{HION-2016-06,HION-2019-03}.
    In this analysis no event-mixing is done to account for acceptance in \Dphi.
  }.
Figure~\ref{fig:example_2pc} shows the distributions for
  central \PbPb\ events (top) panels, 
  peripheral \PbPb\ events (middle panels)
  and \pp\ events bottom panels.
In the \pp\ collisions the Underlying Event (UE) pairs
  are fairly small and most of the pairs are under 
  the away-side peak. 
For the central \PbPb\ events significant number of 
  UE pairs are present, which are themselves modulated in \Dphi\
  due to the effects of collective flow.
These UE pairs must be removed -- both the \PbPb\ and \pp\ measurements -- 
  before the yields and widths on the away-side are evaluated.
The procedure to remove the UE pairs is discussed later.

\begin{figure}[h] 
\centering
\includegraphics[width=1.0\textwidth]%
{figures/AllFigs/can_dphi_fits_cent12_LorentzZYAM_Tight_pTBar_GT5_EffCor}
\includegraphics[width=1.0\textwidth]%
{figures/AllFigs/can_dphi_fits_cent8_LorentzZYAM_Tight_pTBar_GT5_EffCor}
\includegraphics[width=1.0\textwidth]%
{figures/AllFigs/can_dphi_fits_cent11_LorentzZYAM_Tight_pTBar_GT5_pp2017_EffCor}
\caption{
Example two-particle distribution in $\Delta\phi$ for
  muon pairs with |\Deta|>0.8.
The plots correspond to pairs where both muons have 
  the \textit{same-charge} (left), \textit{opposite-charge} (middle)
  and the combined distribution (right).
The lines indicate fits used to extract yields
  on the away-side (discussed later in Section~\ref{sec:DphiFits}).
The top panels correspond to central \PbPb\ events,
  the middle panels correspond to peripheral \PbPb\ events,
  and the bottom panels correspond to \pp\ events.
Note the $y$-ranges of the panels are very different,
  and they are zero suppressed for the \PbPb\ plots.
\label{fig:example_2pc}
}
\end{figure}


Some additional selections are applied on the muon-pairs
  on top of the single-muon selections discussed above.
These are:
\begin{itemize}
\item Cuts on the muons acoplanarity and \pt-asymmetry (defined in Eqns.~\ref{eq:acop}--\ref{eq:asym} below). 
\item Cuts on invariant mass.
\item Cuts on the pseudo-rapidity separation,|\Deta|=|\etaa-\etab|, 
      between the two muons.
\end{itemize}
The Cuts and the reason for applying them are discussed below.
These cuts are applied for both, the \PbPb\ and \pp\ data.

\textbf{Cuts on acoplanarity and \pt-asymmetry}:
These cuts are included to remove muon-pairs that arise from 
  the \gamgammumu\ process~\cite{HION-2018-11,HION-2020-10}, and are only applied to
  pairs where the two muons have opposite-charge.
The acoplanarity and asymmetry are defined as:
\begin{itemize}
  \item Acoplanarity (\Aco): This quantifies how back-to-back in $\Dphi$ the two muons are.
    \begin{eqnarray}
    \Aco \equiv 1- \frac{|\Delta\phi|}{\pi}.
    \label{eq:acop}
    \end{eqnarray}
  \item Asymmetry (\Amumu): This quantifies the \pt\ mismatch between the two muons.
  \begin{eqnarray}
  \Amumu \equiv \frac{|\ptplus-\ptminus|}{\ptplus+\ptminus}.
  \label{eq:asym}
  \end{eqnarray}
\end{itemize}
The \gamgammumu\ process produces dimuons with small values of \Aco\ 
  and \Amumu~\cite{HION-2018-11}.
Thus the following selection is imposed on the opposite-sign pairs 
  to remove the \gamgammumu\ contributions: pairs with $\Aco<0.01$ and $\Amumu<0.08$ are removed.
In principle these cuts should only be applied for the \PbPb\ data, where the \gamgammumu\ process occurs.
However, for consistency and to keep acceptance effects identical 
  between the \pp\ and \PbPb\ analyses, these cuts are applied to the \pp\ data as well.
This cut rejects less than 1\% of opposite-sign pairs in the \pp\ events (where no such background is expected).


\textbf{Invariant mass cuts:} 
Muon-pairs from decays of mesons $\rho$, $\omega$, $\phi$, $J/\psi$, $\psi'$, $\Upsilon$
  etc are excluded by applying cuts on the invariant mass of 
  the pair. 
These cuts are only applied on the pairs where the muons 
  have opposite charge.
The following regions of \Mmumu are removed: 
  0.6--1.05~\GeV,
  2.9--3.3~\GeV\, 
  3.6--3.8~\GeV\, 
  9.2--10.4~\GeV
  and
  70-110~\GeV.
Figure~\ref{fig:minv_cuts} demonstrates the cuts that are applied 
  to reject regions of \Mmumu\ around the masses of the mesons.
the effect of the \Deta\ cuts (discussed below) are also shown.
The \Deta\ cuts remove all pairs below \Mmumu\ of $\sim$3.4~\GeV\ 
 effectively making the mass cuts below this threshold redundant. 
The mass cuts however introduce acceptance effects that are 
  dependent on \Dphi.
Thus a \Dphi\ dependent correction needs to be introduced to 
  the measurement to correct for this.
This correction is discussed in Section~\ref{sec:minv_cut_correction}.

\begin{figure}[h] 
\centering
\includegraphics[width=1.0\textwidth]%
{figures/AllFigs/can_minv2_cent_Tight_pTBar_GT5}
\includegraphics[width=1.0\textwidth]%
{figures/AllFigs/can_minv2_cent_Tight_pTBar_GT5_pp2017}
\caption{
Left panel: \Mmumu\ distributions prior to applying any rejections.
Middle panel: \Mmumu\ distributions after requiring $|\Deta|>0.8$.
Right panel: The distributions after further applying cuts on \Mmumu\ to the
  opposite-charge pairs.
The top and bottom row correspond to \PbPb\ and \pp\ data, respectively.
\label{fig:minv_cuts}
}
\end{figure}

\textbf{\Deta\ cuts}:
A requirement of $|\Deta|>0.8$ is also imposed on the muon-pairs.
The reason to implement this cut is to remove the large number of 
  pairs that are present at $(|\Deta|,\Dphi)\sim(0,0)$.
While the analysis aims to study the muon-pairs on the away-side 
  ($\Dphi=\pi$) that arise from decays of back-to-back HF particles,
  the pair-distribution on the near-side is used to estimate and 
  correct for the UE pairs on the away-side.
Before this is done, the excess pairs at $(|\Deta|,\Dphi)\sim(0,0)$
  need to be removed, and for this reason the $|\Deta|>0.8$ 
  requirement is imposed.
The $|\Deta|>0.8$ cut also removes all resonances below
  a mass of $\sim$3.4~\GeV, as shown in Figure~\ref{fig:minv_cuts}.



Figure~\ref{fig:cutflow} shows the cutflow for the dimuon-pairs
  for the Medium and Tight working points,
  for both the \PbPb\ and \pp\ data.
Figures~\ref{fig:eta_phi_PbPb} and ~\ref{fig:eta_phi_pp}, show
 the $\eta$--$\phi$ distributions for muons from pairs
 that pass all the selections listed above, 
 for the \PbPb\ and \pp\ data, respectively.



\begin{figure}[h] 
\centering
\includegraphics[width=0.5\textwidth]%
{figures/AllFigs/can_CutFlow_pTBar_GT5_EffCor}%
\includegraphics[width=0.5\textwidth]%
{figures/AllFigs/can_CutFlow_Tight_pTBar_GT5_EffCor}
\includegraphics[width=0.5\textwidth]%
{figures/AllFigs/can_CutFlow_pTBar_GT5_pp2017_EffCor}%
\includegraphics[width=0.5\textwidth]%
{figures/AllFigs/can_CutFlow_Tight_pTBar_GT5_pp2017_EffCor}
\caption{
The number of dimuon pairs progressively passing the different
  selection requirements. 
The left and right panels correspond to Medium and Tight 
  muons, respectively.
The top and bottom rows correspond to the \PbPb\ and 
  \pp\ data, respectively.
\label{fig:cutflow}
}
\end{figure}



\begin{figure}[h] 
\centering
\includegraphics[width=1.0\textwidth]%
{figures/AllFigs/can_EtaPhi_pTBar_GT5_EffCor}
\includegraphics[width=1.0\textwidth]%
{figures/AllFigs/can_EtaPhi_Tight_pTBar_GT5_EffCor}
\caption{
$\eta$ and $\phi$ distributions for muons 
  from muon-pairs that pass all the selection requirements.
The top and bottom rows correspond to Medium and 
  Tight muons, respectively.
Plots are for the \PbPb\ data. 
\label{fig:eta_phi_PbPb}
}
\end{figure}


\begin{figure}[h] 
\centering
\includegraphics[width=1.0\textwidth]%
{figures/AllFigs/can_EtaPhi_pTBar_GT5_pp2017_EffCor}
\includegraphics[width=1.0\textwidth]%
{figures/AllFigs/can_EtaPhi_Tight_pTBar_GT5_pp2017_EffCor}
\caption{
Same as Figure~\ref{fig:eta_phi_PbPb}, but for the \pp\ data. 
\label{fig:eta_phi_pp}
}
\end{figure}



\subsubsection{Correction accounting for inefficiency from \Mmumu\ and \gamgammumu\ cuts\label{sec:minv_cut_correction}}
The invariant mass cuts that are applied to the opposite sign-pairs introduce 
  acceptance effects that are $\Dphi$ dependent.
This is demonstrated in Figure~\ref{fig:minv_deta_dphi},
  which shows that the 9.40~\GeV<$\Mmumu$<10.2~\GeV\ mass cut
  which is used to remove the $\Upsilon$ and its excitations
  affects the $\Dphi$ distribution of the pairs.
The \Mmumu cuts below 3.8~\GeV\ are mostly restricted to \Deta<0.8
  and have only a minor effect on the measurement (but are still are corrected for as described below).
Note that while Figure~\ref{fig:minv_deta_dphi} shows the effect of 
  the \Mmumu\ cuts for both same-sign and opposite-sign pairs, 
  in the analysis the are only applied to the opposite-sign pairs.
Thus the correction discussed in this section is relevant only for 
  the opposite-sign pairs (which are then propagated into the All-pairs case).

\begin{figure}[h] 
\centering
\includegraphics[width=1.0\textwidth]%
{figures/AllFigs/can_DetaDphi_MinvCut0_Tight_pTBar_GT5_EffCor}
\caption{
The $\Deta-\Dphi$ distribution of the pairs in various \Mmumu ranges.
The top, middle and bottom rows corresponds to the ranges of 
  0<\Mmumu<125~\GeV (i.e. no cuts on \Mmumu), 9.20~\GeV<\Mmumu<10.4~\GeV.
  and 0<\Mmumu<3.8~\GeV, respectively.
The left, middle and right panels correspond to same-sign pairs
  opposite-sign pairs, and All-pairs, respectively.
Plots are for the \PbPb\ data.
\label{fig:minv_deta_dphi}
}
\end{figure}


Besides the acceptance effects arising from the \Mmumu\ cuts,
  \Dphi\ dependent effects are also introduced because of
  the cuts on the muons acoplanarity and \pt-asymmetry (Eqns.~\ref{eq:acop}--\ref{eq:asym}).
The procedure below also corrects for these cuts.
The $\Dphi$ dependent efficiency effects of the \Mmumu, 
  acoplanarity and \pt-asymmetry cuts on the 
  opposite-sign pairs are estimated by the following two methods:
\begin{enumerate}
  \item By applying the cuts on the ``same-sign'' pairs, and determining
    what fraction of pairs are kept as a function of \Dphi.
  \item By an event-mixing procedure: 
    In this method a single muon triggered (\verb=HLT_mu4=) sample was used
    to select events.
    Pair distributions for muons were build by pairing muons 
      (that were matched to the \verb=HLT_mu4= trigger) from separate events.
    The events from which the muons were combined, were required to be from 
      similar centrality ranges.
    All single-muon and muon-pair cuts that are applied to in the analysis were 
      applied to such ``Event-Mixed''muon-pairs, and the pairs were weighted 
      by the relevant trigger and reconstruction efficiency corrections.
    The effects of the \Mmumu\ and \gamgammumu\ cuts were evaluated by estimating 
      what fraction of the mixed-event pairs survive the 
      cuts as a function of \Dphi.
\end{enumerate}



Figure~\ref{fig:minv_eff} shows the estimated efficiency introduced
  due to the \Mmumu\ cuts as well as the \gamgammumu\ cuts.
To correct for this efficiency the opposite-sign pairs are weighted
  by the inverse of the efficiency as a function of $\Dphi$.
The black points indicate the values estimated using the same-sign
  pairs (item-1 in the list above).
The open blue and red points indicate the efficiency estimated
  using the mixed-events (item-2 in the list above).
The mixed-event efficiency is shows for the same-sign pairs 
  as well as the opposite-sign pairs.
The mixed-event efficiency from the same and opposite-sign
  pairs are virtually identical.
There is however some minor differences between the same-event 
  and mixed-event estimates of the efficiency for $\Dphi>\pi/2$.
This presumably arises as the same-event pairs can have physical correlations
  in \pt\ and $\eta$, which the event-mixing cannot reproduce.
%Also, the event-mixed efficiencies do not show any significant 
%  centrality dependence, while the same-event estimate of the efficieny
%  shows some centrality dependence for $\Dphi>\pi/2$.
%In most central events the same-event and mixed-event estimates of the 
%  efficiency are most similar, and the difference increases for more peripheral events.
%This also arises because in more central events, most of the same-event pairs
%  are combinatoric while in peripheral events the pairs are more likely 
%  to be correlated.
For this reason the efficiency estimated from the same-event distributions
  are used as the default efficiency, with the alternate efficiency obtained from the 
  mixed events used as an estimate systematic uncertainty.


\begin{figure}[h] 
\centering
\includegraphics[width=1.0\textwidth]%
{figures/AllFigs/can_MinvCutEfficiencyFromEventMixing_Tight_pTBar_GT5_EffCor}
\caption{
The estimated efficiency introduced due to the \Mmumu\ cuts, 
  as well as the acoplanarity and asymmetry cuts to remove
  the \gamgammumu\ pairs (see text).
Each panel indicates a different centrality interval.
The last panel showd the efficiency for the 0--100\% centrality bin.
\label{fig:minv_eff}
}
\end{figure}



\subsubsection{Correction for Drell-Yan pairs\label{sec:DY_correction}}
A small contribution to the opposite-sign pairs arises from the Drell-Yan process.
This also needs to be corrected for before extracting the features 
  of the HF pairs.
This correction is done by evaluating the DY contribution using an official
  \powheg\ MC sample described in Ref.~\cite{ATLHI-331}.
This sample was previously used to estimate the DY-background in a previous
  ATLAS measurement~\cite{HION-2020-10} that studied \gamgammumu\ 
  process in heavy-ion collision.
The DY contribution in a given centrality is calculated as:
\begin{align}
	\frac{dN_\mathrm{DY}}{d\Delta\phi} = \mathcal{L}\sigma_\mathrm{had}^\mathrm{Pb+Pb}\times \Delta\mathrm{cent}\times\langle\TAA\rangle\times\frac{d\sigma_\mathrm{DY,NN}}{d\Delta\phi}
\end{align}
where,
  $\sigma_\mathrm{had}^\mathrm{Pb+Pb}$ is the total \PbPb\ hadronic cross-section,
  $\Delta\mathrm{cent}$ represents the width of the centrality interval
    as a fraction (i.e. the width in percentile/100),
  $\langle\TAA\rangle$ is the nuclear-thickness function for the centrality interval (values taken from Ref.~\cite{TAAValues}),
  and $\frac{d\sigma_\mathrm{DY,NN}}{d\Delta\phi}$ is the 
    differential DY cross-section (Ref.~\cite{ATLHI-331}), when applying the same
    fiducial selections as those in the data.

Figure~\ref{fig:DY_Corr_PbPb} shows the pair distributions 
  with/without the DY corrections.
The DY contributons are estimated using three different PDF sets.
In the final measurements, the spread between using the three PDFs 
  is included as a systematic uncertainty, as described later in
  Section~\ref{sec:systematics_drell_yan}.


\begin{figure}[h] 
\centering
\includegraphics[width=1.0\textwidth]%
{figures/AllFigs/can_DY_cent12_Tight_pTBar_GT5_EffCor}
\includegraphics[width=1.0\textwidth]%
{figures/AllFigs/can_DY_cent8_Tight_pTBar_GT5_EffCor}
\includegraphics[width=1.0\textwidth]%
{figures/AllFigs/can_DY_cent11_Tight_pTBar_GT5_pp2017_EffCor}
\caption{
The \Dphi\ distributions with/without the DY corrections.
The DY contributons are estimated using three different PDF sets.
The top and middle rows correspond to the 0--10\% and 60--80\%
  centrality intervals.
The bottom row corresponds to the \pp\ data.
\label{fig:DY_Corr_PbPb}
}
\end{figure}
